\documentclass[11pt]{article}

\usepackage[a4paper,margin=1in]{geometry}
\usepackage{amsmath,amssymb}
\usepackage{enumitem}
\usepackage{titlesec}
\usepackage{booktabs}

\titleformat{\section}{\large\bfseries}{\thesection}{1em}{}

\begin{document}

\begin{center}
{\LARGE \textbf{�� Lecture Scribe: Inequalities and Tail Bounds (L21–L22)}}
\end{center}

\vspace{0.5cm}

\section*{1. Tail of a Random Variable}

\subsection*{Definition}
The \textbf{tail} of a random variable $X$ is:
\[
P(X > x)
\]

\subsection*{Step-by-step Understanding}
\begin{itemize}
\item We focus on \textbf{values greater than a threshold $(x)$}
\item This represents \textbf{extreme events (rare cases)}
\end{itemize}

\subsection*{Why tails matter}
\begin{itemize}
\item Mean $\rightarrow$ gives \textbf{average behavior (easy)}
\item Tail $\rightarrow$ captures \textbf{rare extreme events (hard)}
\end{itemize}

Examples from context:
\begin{itemize}
\item Jobs delayed for long time
\item Server overload (too many requests)
\end{itemize}

\subsection*{Key difficulty}
\begin{itemize}
\item Computing exact tail probability requires:
\begin{itemize}
\item Summing many small probabilities
\end{itemize}
\item Hence $\rightarrow$ \textbf{complex and computationally expensive}
\end{itemize}

\section*{2. Why Tail Bounds?}

\subsection*{Problem}
Exact computation:
\[
P(X \ge k)
\]
involves long summations (Binomial/Poisson).

\subsection*{Solution Idea}
Instead of exact value:
\begin{itemize}
\item Find an \textbf{upper bound}
\item Easier to compute
\item Still guarantees safety
\end{itemize}

\subsection*{Key Insight}
\[
P(X \ge k) \le \text{simple bound}
\]

\section*{3. Markov’s Inequality}

\subsection*{Definition}
For a \textbf{non-negative random variable $(X)$}:
\[
P(X \ge a) \le \frac{E[X]}{a}
\]

\subsection*{Assumptions}
\begin{itemize}
\item $X \ge 0$
\item Only \textbf{mean $(E[X])$} is known
\end{itemize}

\subsection*{Proof Idea (Step-by-step)}
\begin{enumerate}
\item Start with expectation:
\[
E[X] = \sum x \cdot P(X=x)
\]
\item Focus only on values where $X \ge a$
\item Replace those values by $a$:
\[
E[X] \ge a \cdot P(X \ge a)
\]
\item Rearranging:
\[
P(X \ge a) \le \frac{E[X]}{a}
\]
\end{enumerate}

\subsection*{Interpretation}
\begin{itemize}
\item Large values are rare
\item Uses \textbf{only mean $\rightarrow$ very weak bound}
\end{itemize}

\subsection*{Solved Example (Step-by-step)}

\textbf{Given:}
\begin{itemize}
\item Average grade = 70
\item Threshold = 90
\end{itemize}

\textbf{Step 1: Define variable}
\[
X = \text{student grade}, \quad E[X] = 70
\]

\textbf{Step 2: Apply Markov}
\[
P(X \ge 90) \le \frac{70}{90}
\]

\textbf{Step 3: Simplify}
\[
P(X \ge 90) \le \frac{7}{9} \approx 0.778
\]

\textbf{Result}
\begin{itemize}
\item At most \textbf{77.8\% students} can get A
\end{itemize}

\section*{4. Chebyshev’s Inequality}

\subsection*{Key Idea}
Uses:
\begin{itemize}
\item Mean $(\mu)$
\item Variance $(\sigma^2)$
\end{itemize}

Controls:
\begin{itemize}
\item \textbf{Deviation from mean}
\end{itemize}

\subsection*{Definition}
\[
P(|X - \mu| \ge a) \le \frac{\text{Var}(X)}{a^2}
\]

\subsection*{Proof Idea}
\begin{enumerate}
\item Start with deviation:
\[
|X - \mu|
\]
\item Square it:
\[
(X - \mu)^2
\]
\item Apply Markov inequality:
\[
P((X-\mu)^2 \ge a^2)
\le \frac{E[(X-\mu)^2]}{a^2}
\]
\item Use:
\[
E[(X-\mu)^2] = \text{Var}(X)
\]
\end{enumerate}

\subsection*{Interpretation}
\begin{itemize}
\item Bounds \textbf{how far values go from mean}
\item Better than Markov
\end{itemize}

\subsection*{Solved Example}

\textbf{Given:}
\begin{itemize}
\item Mean = 70
\item Std dev = 5 $\rightarrow$ Variance = 25
\item Threshold = 90
\end{itemize}

\textbf{Step 1: Convert condition}
\[
X \ge 90 \Rightarrow X - 70 \ge 20
\]

\textbf{Step 2: Apply Chebyshev}
\[
P(X \ge 90) \le P(|X - 70| \ge 20)
\]

\textbf{Step 3: Substitute}
\[
\le \frac{25}{20^2} = \frac{25}{400}
\]

\textbf{Step 4: Simplify}
\[
= 0.0625
\]

\textbf{Result}
\begin{itemize}
\item Only \textbf{6.25\% maximum probability}
\end{itemize}

\section*{5. Chernoff Bound}

\subsection*{Key Idea}
Progression:
\begin{itemize}
\item Markov $\rightarrow$ uses $(X)$
\item Chebyshev $\rightarrow$ uses $((X-\mu)^2)$
\item Chernoff $\rightarrow$ uses $(e^{tX})$
\end{itemize}

\subsection*{Insight}
\begin{itemize}
\item Exponential grows fast
\item Amplifies large values $\rightarrow$ tighter tail bound
\end{itemize}

\subsection*{Core Trick}
\[
P(X \ge a) = P(e^{tX} \ge e^{ta})
\]

Apply Markov:
\[
P(X \ge a) \le \frac{E[e^{tX}]}{e^{ta}}
\]

\subsection*{Final Form}
\[
P(X \ge a) \le \min_{t>0} \frac{E[e^{tX}]}{e^{ta}}
\]

\subsection*{Key Insight}
\begin{itemize}
\item Holds for all $t > 0$
\item Choose best $t \rightarrow$ tightest bound
\end{itemize}

\subsection*{Lower Tail Bound (Derivation Steps)}

Goal:
\[
P(X \le a)
\]

\textbf{Step 1: Choose $t < 0$}
\begin{itemize}
\item Multiplying flips inequality
\end{itemize}

\[
P(X \le a) = P(tX \ge ta)
\]

\textbf{Step 2: Exponentiate}
\[
P(e^{tX} \ge e^{ta})
\]

\textbf{Step 3: Apply Markov}
\[
\le \frac{E[e^{tX}]}{e^{ta}}
\]

\textbf{Step 4: Minimize over $t < 0$}

Final:
\[
P(X \le a) \le \min_{t<0} e^{-ta} E[e^{tX}]
\]

\subsection*{Example (Binomial Case)}

Given:
\begin{itemize}
\item $X \sim \text{Binomial}(20, 0.5)$
\item $\mu = 10$
\end{itemize}

Goal:
\[
P(X = 0)
\]

\textbf{Step 1: Use lower tail}
\[
(1 - \delta)\mu = 0 \Rightarrow \delta = 1
\]

\textbf{Step 2: Apply bound}
\[
P(X \le 0) \le e^{-\mu \delta^2 / 2}
\]

\textbf{Step 3: Substitute}
\[
= e^{-10/2} = e^{-5}
\]

\textbf{Result}
\[
\approx 0.0067
\]

\begin{itemize}
\item Very small probability
\end{itemize}

\section*{6. Jensen’s Inequality}

\subsection*{Convex Function Definition}
A function $g(x)$ is convex if:
\[
g(\alpha x_1 + (1-\alpha)x_2)
\le \alpha g(x_1) + (1-\alpha)g(x_2)
\]

\subsection*{Key Interpretation}
\begin{itemize}
\item Function lies \textbf{below straight line}
\item Convex combination $\rightarrow$ weighted average
\end{itemize}

\subsection*{From Convexity to Expectation}
Expectation = weighted average:
\[
E[X] = \sum p_i x_i
\]

Apply convexity:
\[
g(E[X]) \le E[g(X)]
\]

\subsection*{Jensen’s Inequality (Final)}
\[
g(E[X]) \le E[g(X)]
\]

\subsection*{Interpretation}
\begin{itemize}
\item Apply average first $\rightarrow$ smaller value
\item Apply function first $\rightarrow$ larger value
\end{itemize}

\subsection*{Example (Step-by-step)}

Given:
\begin{itemize}
\item $X = 1$ with 0.5
\item $X = 3$ with 0.5
\item $g(x) = x^2$
\end{itemize}

\textbf{Step 1: Compute mean}
\[
E[X] = 0.5(1) + 0.5(3) = 2
\]

\textbf{Step 2: Left side}
\[
g(E[X]) = 2^2 = 4
\]

\textbf{Step 3: Right side}
\[
E[g(X)] = 0.5(1^2) + 0.5(3^2) = 5
\]

\textbf{Result}
\[
4 \le 5
\]

\section*{7. Case Study: System Design Insight}

\subsection*{Problem Setup}
\begin{itemize}
\item $X$: incoming requests/sec
\item $C$: system capacity
\end{itemize}

Failure:
\[
X > C
\]

Requirement:
\[
P(X > C) \le 0.01
\]

\subsection*{Results Using Different Bounds}

\begin{center}
\begin{tabular}{lll}
\toprule
Method & Info Used & Capacity Required \\
\midrule
Markov & Mean only & 1,000,000 \\
Chebyshev & Mean + variance & 30,000 \\
Chernoff & Full distribution & 14,500 \\
\bottomrule
\end{tabular}
\end{center}

\subsection*{Final Insight}

\textbf{Golden Rule:}

\begin{quote}
More statistical information $\Rightarrow$ tighter bounds $\Rightarrow$ lower cost
\end{quote}

\section*{Final Summary (Exam Ready)}

\begin{itemize}
\item Tail = probability of extreme events
\item Markov:
\begin{itemize}
\item Uses mean
\item Weak bound
\end{itemize}
\item Chebyshev:
\begin{itemize}
\item Uses mean + variance
\item Better bound
\end{itemize}
\item Chernoff:
\begin{itemize}
\item Uses exponential transformation
\item Tightest bound
\end{itemize}
\item Jensen:
\begin{itemize}
\item Convexity $\Rightarrow$ expectation inequality
\end{itemize}
\end{itemize}

\end{document}